% Part: first-order-logic
% Chapter: completeness
% Section: identity

\documentclass[../../../include/open-logic-section]{subfiles}

\begin{document}

\olfileid[hi]{fol}{com}{ide}
\olsection{तादात्म्य}

\begin{explain}
पिछले अनुभाग में दिया पद-प्रतिरूप का निर्माण उन समुच्चयों~$\Gamma$ के लिए
प्रथम-क्रम तर्क की पूर्णता स्थापित करने को पर्याप्त है जिनमें~$\eq$ नहीं
है। पद-प्रतिरूप प्रत्येक ऐसे $!A \in \Gamma^*$ को संतुष्ट करता है जिसमें
~$\eq$ नहीं है (और इसलिए सभी $!A \in \Gamma$ को भी)। किंतु यदि $\eq$
उपस्थित हो तो यह काम नहीं करता। कारण यह है कि तब $\Gamma^*$ में कोई
!!a{sentence}~$\eq[t][t']$ हो सकता है, जबकि पद-प्रतिरूप में किसी पद का मान
स्वयं वही पद होता है। अतः यदि $t$ और $t'$ भिन्न पद हैं, तो पद-प्रतिरूप में
उनके मान---क्रमशः $t$ और $t'$---भिन्न हैं, इसलिए $\eq[t][t']$ असत्य है।
फिर भी हम इसे ``गुणनखंडन'' नामक रचना से ठीक कर सकते हैं।
\end{explain}

\begin{defn}
  $\Gamma^*$ को~$\Lang L$ में वाक्यों का संगत और !!{complete} समुच्चय मानें।
  हम संबंध $\approx$ को~$\Lang L$ के बंद पदों के समुच्चय पर इस प्रकार
  परिभाषित करते हैं:
  \[
  t \approx t' \text{\quad तभी और केवल तभी \quad} \eq[t][t'] \in \Gamma^*
  \]
\end{defn}

\begin{prop}
\ollabel{prop:approx-equiv}
संबंध $\approx$ में निम्न गुण हैं:
\begin{enumerate}
\item $\approx$ स्वतुल्य है।
\item $\approx$ सममित है।
\item  $\approx$ संक्रामक है।
\item यदि $t \approx t'$, $f$ कोई !!a{function} है, और $t_1$, \dots,
  $t_{i-1}$, $t_{i+1}$, \dots, $t_n$ बंद पद हैं, तो
\[
\Atom{f}{t_1,\dots, t_{i-1}, t, t_{i+1}, \dots, t_n} \approx
\Atom{f}{t_1,\dots, t_{i-1}, t', t_{i+1}, \dots, t_n}.
\]
\item यदि $t \approx t'$, $R$ कोई !!a{predicate} है, और $t_1$, \dots,
  $t_{i-1}$, $t_{i+1}$, \dots, $t_n$ बंद पद हैं, तो
\begin{multline*}
\Atom{R}{t_1,\dots, t_{i-1}, t, t_{i+1}, \dots, t_n} \in \Gamma^* \text{ तभी और केवल तभी } \\
\Atom{R}{t_1,\dots, t_{i-1}, t', t_{i+1}, \dots, t_n} \in \Gamma^*.
\end{multline*}
\end{enumerate}
\end{prop}

\begin{proof}
चूँकि $\Gamma^*$ संगत और !!{complete} है, $\eq[t][t'] \in
\Gamma^*$ तभी और केवल तभी जब $\Gamma^* \Proves \eq[t][t']$। अतः निम्न
कथनों को दिखाना पर्याप्त है:
\begin{enumerate}
\item $\Gamma^* \Proves \eq[t][t]$, सभी बंद पदों~$t$ के लिए।
\item यदि $\Gamma^* \Proves \eq[t][t']$, तो $\Gamma^* \Proves \eq[t'][t]$।
\item यदि $\Gamma^* \Proves \eq[t][t']$ और $\Gamma^* \Proves
  \eq[t'][t'']$, तो $\Gamma^* \Proves \eq[t][t'']$।
\item यदि $\Gamma^* \Proves \eq[t][t']$, तो
\[
\Gamma^* \Proves
\eq[\Atom{f}{t_1,\dots,t_{i-1},t,t_{i+1},,\dots,t_n}][\Atom{f}{t_1,\dots,t_{i-1},t',t_{i+1},\dots,t_n}]
\]
प्रत्येक $n$-स्थानीय !!{function}~$f$ और बंद पदों $t_1$, \dots,
$t_{i-1}$, $t_{i+1}$, \dots,~$t_n$ के लिए।
\item यदि $\Gamma^* \Proves \eq[t][t']$ और
$\Gamma^* \Proves
\Atom{R}{t_1,\dots,t_{i-1},t,t_{i+1},\dots,t_n}$, तो
$\Gamma^* \Proves \Atom{R}{t_1,\dots,t_{i-1},t',t_{i+1},\dots,t_n}$
प्रत्येक $n$-स्थानीय !!{predicate}~$R$ और बंद पदों $t_1$, \dots,
$t_{i-1}$, $t_{i+1}$, \dots,~$t_n$ के लिए।
\end{enumerate}
\end{proof}

\begin{prob}
\olref[fol][com][ide]{prop:approx-equiv} का प्रमाण पूरा कीजिए।
\end{prob}

\begin{defn}
मान लें कि $\Gamma^*$, भाषा~$\Lang L$ में संगत और !!{complete} समुच्चय है,
$t$ बंद पद है, और $\approx$ पिछली परिभाषा जैसा है। तब:
\[
\equivrep{t}{\approx} = \Setabs{t'}{t'\in \Trm[L], t \approx t'}
\]
और $\equivclass{\Trm[L]}{\approx} = \Setabs{\equivrep{t}{\approx}}{t \in \Trm[L]}$।
\end{defn}

\begin{defn}
\ollabel{defn:term-model-factor}
मान लें कि $\Struct M = \Struct M(\Gamma^*)$,~$\Gamma^*$ के लिए
\olref[mod]{defn:termmodel} वाला पद-प्रतिरूप है। तब
$\Struct{\equivclass{M}{\approx}}$ निम्न !!{structure} है:
\begin{enumerate}
\item $\Domain{\equivclass{M}{\approx}} = \equivclass{\Trm[L]}{\approx}$।
\item $\Assign{c}{\equivclass{M}{\approx}} = \equivrep{c}{\approx}$
\item $\Assign{f}{\equivclass{M}{\approx}}(\equivrep{t_1}{\approx}, \dots,
  \equivrep{t_n}{\approx}) = \equivrep{\Atom{f}{t_1,\dots, t_n}}{\approx}$
\item $\tuple{\equivrep{t_1}{\approx}, \dots, \equivrep{t_n}{\approx}} \in
  \Assign{R}{\equivclass{M}{\approx}}$ तभी और केवल तभी जब
  $\Sat{M}{\Atom{R}{t_1,\dots, t_n}}$, अर्थात् तभी और केवल तभी जब
  $\Atom{R}{t_1,\dots,
  t_n} \in \Gamma^*$।
\end{enumerate}
\end{defn}

\begin{explain}
ध्यान दें कि हमने $\Assign{f}{\equivclass{M}{\approx}}$ और
$\Assign{R}{\equivclass{M}{\approx}}$ को~$\equivclass{\Trm[L]}{\approx}$
के अवयवों के लिए, उन्हें $\equivrep{t}{\approx}$
कहकर, अर्थात् \emph{प्रतिनिधियों}~$t \in
\equivrep{t}{\approx}$ के माध्यम से परिभाषित किया है। हमें सुनिश्चित करना
होगा कि ये परिभाषाएँ इन प्रतिनिधियों के चुनाव पर निर्भर न हों; अर्थात् ऐसे
किसी दूसरे चुनाव~$t'$ के लिए, जो वही तुल्यता-वर्ग निर्धारित करता है
($\equivrep{t}{\approx} = \equivrep{t'}{\approx}$), परिभाषाएँ वही परिणाम
दें। उदाहरणार्थ, यदि $R$ एक-स्थानीय !!{predicate} है, तो परिभाषा का अंतिम
खंड कहता है कि $\equivrep{t}{\approx} \in \Assign{R}{\equivclass{M}{\approx}}$
तभी और केवल तभी जब
$\Sat{M}{\Atom{R}{t}}$। यदि किसी दूसरे पद~$t'$ के लिए, जहाँ $t \approx
t'$, $\Sat/{M}{\Atom{R}{t}}$, तो परिभाषा माँगेगी कि
$\equivrep{t'}{\approx} \notin \Assign{R}{\equivclass{M}{\approx}}$। यदि
$t \approx t'$, तो $\equivrep{t}{\approx} =
\equivrep{t'}{\approx}$, पर हमारे पास $\equivrep{t}{\approx}
\in \Assign{R}{\equivclass{M}{\approx}}$ और $\equivrep{t}{\approx}
\notin \Assign{R}{\equivclass{M}{\approx}}$ दोनों नहीं हो सकते। किंतु
\olref{prop:approx-equiv} गारंटी देता है कि ऐसा नहीं हो सकता।
\end{explain}

\begin{prop}
$\Struct{\equivclass{M}{\approx}}$ सुनिर्धारित है; अर्थात् यदि $t_1$,
  \dots, $t_n$, $t_1'$, \dots, $t_n'$ बंद पद हैं,
  और $t_i \approx t_i'$ तो
\begin{enumerate}
\item $\equivrep{\Atom{f}{t_1,\dots, t_n}}{\approx} =
    \equivrep{\Atom{f}{t_1',\dots, t_n'}}{\approx}$, अर्थात्
  \[
  \Atom{f}{t_1,\dots, t_n} \approx \Atom{f}{t_1',\dots, t_n'}
  \]
  और
\item $\Sat{M}{\Atom{R}{t_1,\dots, t_n}}$ तभी और केवल तभी जब
  $\Sat{M}{\Atom{R}{t_1',\dots, t_n'}}$, अर्थात्
  \[
    \Atom{R}{t_1,\dots, t_n} \in \Gamma^* \text{ तभी और केवल तभी }
    \Atom{R}{t_1',\dots, t_n'} \in \Gamma^*.
  \]
\end{enumerate}
\end{prop}

\begin{proof}
\olref{prop:approx-equiv} से~$n$ पर आगमन द्वारा निकलता है।
\end{proof}

पद-प्रतिरूप की तरह ही, सत्यता उपपत्ति सिद्ध करने से पहले हमें निम्न उपपत्ति
चाहिए।

\begin{lem} 
\ollabel{lem:val-in-termmodel-factored} मान लें कि $\Struct M = \Struct M
 (\Gamma^*)$; तब $\Value{t}{\equivclass{M}{\approx}} = \equivrep{t}
 {\approx}$।
\end{lem}
\begin{proof} 
प्रमाण \olref[mod]{lem:val-in-termmodel} के प्रमाण के समान है।
\end{proof}

\begin{prob}
\olref[fol][com][ide]{lem:val-in-termmodel-factored} का प्रमाण पूरा कीजिए।
\end{prob}

\begin{lem}
\ollabel{lem:truth}
$\Sat{\equivclass{M}{\approx}}{!A}$ तभी और केवल तभी जब $!A \in \Gamma^*$,
सभी वाक्यों~$!A$ के लिए।
\end{lem}

\begin{proof}
~$!A$ पर आगमन से, ठीक \olref[mod]{lem:truth} के प्रमाण की तरह। केवल उस
स्थिति पर अतिरिक्त ध्यान चाहिए जिसमें $!A \ident
\eq[t][t']$।
\begin{align*}
\Sat{\equivclass{M}{\approx}}{\eq[t][t']} & \text{ तभी और केवल तभी } \equivrep{t}{\approx} = \equivrep{t'}{\approx}
\text{ ($\Struct{\equivclass{M}{\approx}}$ की परिभाषा से)}\\
& \text{ तभी और केवल तभी } t \approx t' \text{ ($\equivrep{t}{\approx}$ की परिभाषा से)}\\
& \text{ तभी और केवल तभी } \eq[t][t'] \in \Gamma^* \text{ ($\approx$ की परिभाषा से)।}
\end{align*}
\end{proof}

\begin{digress}
ध्यान दें कि $\Struct{M(\Gamma^*)}$ सदा !!{enumerable} और अनंत है, किंतु
$\Struct{\equivclass{M}{\approx}}$ परिमित हो सकता है, क्योंकि संभव है कि
केवल परिमिततः अनेक वर्ग $\equivrep{t}{\approx}$ हों। यही अपेक्षित है,
क्योंकि $\Gamma$ में ऐसे !!{sentence}s हो सकते हैं जो अपने को सत्य बनाने
वाली प्रत्येक !!{structure} को परिमित होने को बाध्य करें। उदाहरणार्थ,
$\lforall[x][\lforall[y][x =
y]]$ एक संगत !!{sentence} है, पर इसे संतुष्ट करने वाली !!{structure}s केवल
वे हैं जिनके !!a{domain} में ठीक एक !!{element} है।
\end{digress}

\end{document}
